\documentclass[11pt,a4paper]{article}
\usepackage[utf8]{inputenc}
\usepackage[T1]{fontenc}
\usepackage[spanish]{babel}
\usepackage[margin=23mm]{geometry}
\usepackage{lmodern,xcolor,hyperref,pdfpages}
\definecolor{Navy}{HTML}{102A43}
\definecolor{Blue}{HTML}{2563EB}
\hypersetup{colorlinks=true,urlcolor=Blue,pdfauthor={Jason Chambe},pdftitle={Publicación de LinkedIn: tamaño de muestra en R}}
\setlength{\parindent}{0pt}\setlength{\parskip}{9pt}
\begin{document}
{\small\color{Blue}BIOESTADÍSTICA / TAREA 2}\par
{\LARGE\bfseries\color{Navy}Publicación de LinkedIn}\par
{\large Jason Chambe}\par
{\small Texto propuesto y carrusel de 20 imágenes. Este archivo reúne el material preparado para compartir; no acredita una publicación en la plataforma.}\par
\vspace{4mm}\hrule\vspace{4mm}
¿Cuántos participantes necesita un estudio? Antes de calcular, necesitamos saber qué estamos comparando.

Como parte de mi curso de Bioestadística, preparé un tutorial de tamaño de muestra en R con tres ejercicios: hemoglobina en Puno, anemia en Juliaca y comparación de prevalencias entre dos distritos.

En cada ejercicio desarrollé tres variantes para entender qué ocurre cuando cambia la hipótesis, la potencia o la diferencia que esperamos detectar. Una de las comparaciones me ayudó a verlo con claridad: al reducir de 20 a 10 puntos porcentuales la disminución de anemia que queremos detectar, la muestra pasa de 41 a 178 niños, bajo los mismos supuestos del cálculo bilateral.

El tutorial incluye la justificación del método, el código en R y la interpretación de cada resultado. También comparto una práctica ampliada con ocho escenarios y 24 variantes, en QMD y PDF.

Mi principal aprendizaje: el número que devuelve R depende del diseño y de los supuestos que escribimos. Entenderlos es parte del cálculo.

Puedes descargar el material aquí:\par\url{https://github.com/chambejason951-collab/BIOESTADISTICA/tree/main/TAREA2}

Material elaborado a partir de las clases y la práctica de Antonio M. Quispe, con consulta de la documentación de R y sus paquetes. Los ejemplos son didácticos y explicitan sus supuestos.

¿Qué parte del cálculo de tamaño de muestra te resulta más difícil de justificar?

\#Bioestadística \#RStats \#TamañoDeMuestra \#Aprendizaje
\clearpage
\includepdf[pages=-,fitpaper=true,pagecommand={}]{linkedin/carrusel_tamano_muestra.pdf}
\end{document}
